% TWO PAGES, AND NO MORE. How to keep to the limit, in the order to try things:
%   1. Cut words. Every paragraph that does not support the main claim goes.
%   2. Make the figure smaller: change the width in the \includegraphics / tikzpicture
%      (a figure at one column width is plenty), or let it span less height.
%   3. Cite fewer, shorter references; the list is already set in \small with no space between entries.
%   4. Last resort: \looseness=-1 before a paragraph that ends with a few words on
%      its last line pulls that line back (put it right before the paragraph's end).
% This sample fills about a page and a half, so you have room to write.
% Do NOT shrink the font below 10 pt or the margins below 1.5 cm: reviewers notice, and
% many venues reject it. Check the page count at the end of the build: this template
% prints a package warning in the log as soon as the document runs past two pages.
\documentclass[10pt,a4paper,twocolumn]{article}
\usepackage[T1]{fontenc}
\usepackage[utf8]{inputenc}
\usepackage{newtxtext,newtxmath}    % Times-like text and matching maths: compact and familiar
\usepackage[top=1.8cm,bottom=1.8cm,left=1.7cm,right=1.7cm,columnsep=0.7cm]{geometry}
\usepackage{microtype}              % tighter, better line breaking, which saves lines
\usepackage{amsmath}
\usepackage{graphicx}
\usepackage{booktabs}
\usepackage{pgfplots}               % the figure is drawn from data, so no image file is needed
\pgfplotsset{compat=1.16}
\usepackage[font=small,labelfont=bf,skip=3pt,hypcap=false]{caption}  % also provides \captionof
\usepackage[numbers,sort&compress]{natbib}
\bibliographystyle{abbrvnat}
\setlength{\bibsep}{0pt}            % no space between references
\usepackage[hidelinks]{hyperref}

% Compact headings and paragraphs: the part that makes two pages possible.
\usepackage[compact]{titlesec}
\titleformat{\section}{\normalfont\large\bfseries}{\thesection}{0.6em}{}
\titleformat{\subsection}{\normalfont\normalsize\bfseries}{\thesubsection}{0.5em}{}
\titlespacing*{\section}{0pt}{1.4ex plus .3ex}{0.7ex}
\titlespacing*{\subsection}{0pt}{1.2ex plus .2ex}{0.5ex}
\setlength{\parskip}{0pt}
\setlength{\parindent}{1em}
\setlength{\textfloatsep}{8pt plus 2pt minus 2pt}
\setlength{\floatsep}{6pt plus 2pt minus 2pt}
\pagestyle{plain}

% Warn when the document is longer than two pages (after the last page, the page counter is its number).
\AddToHook{enddocument/afterlastpage}{%
  \ifnum\value{page}>2 \GenericWarning{}{Package warning: this extended abstract is over two pages}\fi}

\begin{document}

% The title block spans both columns: \twocolumn[...] takes it as its optional argument.
\twocolumn[{%
  \begin{center}
    {\Large\bfseries Title of the Extended Abstract: Short and Specific\par}
    \vspace{0.8ex}
    First Author$^{1}$, Second Author$^{1,2}$, Third Author$^{2}$\par
    \vspace{0.4ex}
    {\small $^{1}$Example University, City, Country \quad
     $^{2}$Example Research Institute, City, Country\par
     \texttt{first.author@example.org}\par}
  \end{center}
  \begin{abstract}
  \small\noindent
  Two or three sentences only: the problem, the approach and the main result. An extended
  abstract usually has no separate abstract at all, so check the call for papers; if the venue
  asks for one, keep it this short, or delete this block.
  \end{abstract}
  \vspace{1ex}
  \begin{center}
  % A figure inside the \twocolumn[...] argument spans both columns and stays on page 1.
  % (A float of type figure* would be pushed to page 2.) It cannot float, so keep it small.
  \begin{minipage}[t]{0.48\textwidth}\centering
  \begin{tikzpicture}
    \begin{axis}[ybar, width=\linewidth, height=4.2cm, bar width=11pt, ymin=0, ymax=16,
                 ylabel={Error (\%)}, symbolic x coords={Full,No A,No B,No C},
                 xtick=data, enlarge x limits=0.15, nodes near coords,
                 nodes near coords style={font=\scriptsize},
                 tick label style={font=\footnotesize}, label style={font=\footnotesize}]
      \addplot[fill=blue!40!gray, draw=none] coordinates {(Full,8.1) (No A,9.4) (No B,12.1) (No C,8.9)};
    \end{axis}
  \end{tikzpicture}
  \captionof{figure}{Removing one component at a time.}\label{fig:bars}
  \end{minipage}\hfill
  \begin{minipage}[t]{0.48\textwidth}\centering
  \begin{tikzpicture}
    \begin{axis}[width=\linewidth, height=4.2cm, xlabel={Model size (millions of parameters)},
                 ylabel={Error (\%)}, xmode=log, ymin=0, grid=major, grid style={gray!25},
                 tick label style={font=\footnotesize}, label style={font=\footnotesize}]
      \addplot[thick, mark=*, blue!60!black] coordinates {(1,16) (4,12.5) (16,10.1) (64,8.7) (256,8.1)};
    \end{axis}
  \end{tikzpicture}
  \captionof{figure}{Error against model size.}\label{fig:size}
  \end{minipage}
\end{center}
}]

\section{Introduction}
Say what problem you address and why a reader at this venue should care, in a few
sentences. Cite only the closest work \citep{he2016,vaswani2017}. State your
contribution in one clear sentence: we propose X, which improves Y by Z. Then list
what the reader gets, in at most three points:
\begin{itemize}\itemsep=0pt\topsep=2pt\parsep=0pt
  \item a method that does something concrete and new;
  \item a result that can be checked on a standard benchmark;
  \item an open implementation, with a link (put it in a footnote).
\end{itemize}


\section{Background}
Two short paragraphs of context are enough: what is known, and the gap you fill. A
reviewer reading forty abstracts wants the gap in the first half page. Name one or two
alternatives and say in a clause why they fall short, rather than surveying the field.

\section{Related work}
Name the two or three closest papers and say how yours differs from each in one sentence;
longer surveys do not fit. A reader should know after this paragraph why your work is not
a rerun of \citet{he2016} or an application of \citet{vaswani2017}.

\section{Approach}\label{sec:approach}
Describe the idea with just enough detail to judge it. One equation is usually enough:
\begin{equation}
  \hat{\theta} = \arg\min_{\theta}\; \frac{1}{n}\sum_{i=1}^{n} \ell\bigl(f_{\theta}(x_i),\, y_i\bigr).
  \label{eq:objective}
\end{equation}
Details that a full paper would give, such as proofs and parameter settings, are left
for the full version; say so in one sentence and point to it. The model is trained with
Adam \citep{kingma2015}.

\subsection{Training procedure}
If the method is a procedure, a numbered list is shorter than pseudocode and just as clear:
\begin{enumerate}\itemsep=0pt\topsep=2pt\parsep=0pt
  \item sample a mini-batch and compute the loss in Equation~\eqref{eq:objective};
  \item update the parameters with a step of size $\eta_t$, where
        $\eta_t = \eta_0 / \sqrt{t}$;
  \item every $k$ steps, evaluate on the validation set and keep the best model.
\end{enumerate}
State the one design choice that is new. Everything standard (the optimiser, the schedule,
the initialisation) can be named and cited rather than explained, which is where most of
the space is saved.

\section{Results}\label{sec:results}
Show the one result that carries the argument, and report it honestly: mean over several
runs, the same data split for every method, and the settings used. Reviewers of a two-page
text cannot ask for the appendix, so the claim has to stand on what is on the page. Figure~\ref{fig:curve} is drawn directly
from data in \LaTeX{}; at one column wide it takes about a sixth of a page.

\begin{figure}[t]
  \centering
  \begin{tikzpicture}
    \begin{axis}[width=\columnwidth, height=4.4cm, xlabel={Training steps (thousands)},
                 ylabel={Error (\%)}, xmin=0, xmax=50, ymin=0,
                 legend pos=north east, legend style={font=\footnotesize},
                 tick label style={font=\footnotesize}, label style={font=\footnotesize}]
      \addplot[thick, dashed] coordinates {(0,38) (5,22) (10,17) (20,14.2) (30,13.1) (40,12.7) (50,12.4)};
      \addplot[thick, black] coordinates {(0,36) (5,18) (10,12.5) (20,9.6) (30,8.6) (40,8.2) (50,8.1)};
      \legend{Baseline, Proposed}
    \end{axis}
  \end{tikzpicture}
  \caption{Validation error during training. Make the caption say what to see in the figure.}
  \label{fig:curve}
\end{figure}

\subsection{Setup}
Report the data, the metric and the number of runs in two lines. Everything a reviewer needs
to judge fairness (same budget, same split, tuned baselines) belongs here, because there is
no appendix to hide it in.

\subsection{Comparison}
A small table can replace a paragraph of numbers: Table~\ref{tab:main}.

\begin{table}[t]
  \centering
  \caption{Test error (\%), mean of five runs.}
  \label{tab:main}
  \begin{tabular}{lcc}
    \toprule
    Method & Error (\%) & Time (ms) \\
    \midrule
    Baseline & 12.4 & 38 \\
    Proposed & \textbf{8.1} & 41 \\
    \bottomrule
  \end{tabular}
\end{table}

\subsection{Ablation}
A second experiment, only if it answers a question the first one raises. Figures~\ref{fig:bars}
and~\ref{fig:size} (top of page 1) show which component of the method contributes most and how the result scales.


\section{Discussion}
Say what the result means and where it stops holding. A sentence on limitations makes the
rest more believable. If the venue values it, add a sentence on how the work could be used
or what you would like to discuss with the audience; many workshops select abstracts for
exactly that reason.

\section{Limitations}
Be specific about where the method fails: the data it was not tested on, the assumption it
relies on, the cost it adds. Reviewers trust an abstract that says this itself, and it saves you from the question
in the discussion round. If a limitation is also an opportunity, say that too, in the same sentence.

\section{Outlook}
One short paragraph on the next step: the experiment you will run, the data you are
collecting, or the question you hope the audience can help with. At a workshop this is often
the part people remember, because it invites conversation.

\section{Conclusion}
State what was shown and the next step in two sentences. Do not repeat the introduction.

\section*{Acknowledgements}
Funding and thanks, in one line (delete if your venue counts them against the limit).

{\small
\bibliography{refs}}

\end{document}
